\section{Results III: Explainability as a Test}
\label{sec:xai}

Displaying a few heat maps and asserting that a model ``focuses on clinically
relevant regions'' is unfalsifiable as stated. This section instead states two
predictions in advance and tests each with the tool suited to it.

\textbf{Two tools, two axes.} Our explainability evidence comes from two methods
that describe different things, and keeping them apart is essential to reading
the results. SHAP is applied to the \emph{classical} models (SVM and random
forest, on MFCC and PWP features) and answers a question about \emph{frequency}:
which spectral regions carry the attribution. Grad-CAM is applied to the
\emph{CNN} (on the log-Mel spectrogram) and answers a question about
\emph{time}: where in the cardiac cycle the attribution falls. The two are
orthogonal by construction, and deliberately so. SHAP acts on pooled feature
vectors, from which the time axis has already been discarded, so it can speak
only to frequency; the trustworthiness of that frequency mapping is why it is
used there. Grad-CAM produces a two-dimensional time--frequency map, but its
last convolutional feature map is only a few cells wide, so its frequency axis
is too coarse for spectral claims while its time axis (188 frames) is
fine-grained. We therefore read each method on the axis where it is reliable and
on no other. There is consequently no common axis on which a SHAP frequency
profile and a Grad-CAM temporal profile could be compared directly, and we do
not attempt such a comparison; the cross-method agreement we \emph{can} measure
is agreement among independent SHAP explainers, reported in
Section~\ref{sec:xai-freq}.

\begin{table}[t]
\centering
\caption{What each explainability experiment tests. ``Target'' is the axis the
method is read on; the other axis is not claimed.}
\label{tab:xai-map}
\setlength{\tabcolsep}{4pt}
\begin{tabular}{llll}
\toprule
Experiment & Model & Method & Target \\
\midrule
Weight randomisation & CNN & Grad-CAM & (validity) \\
Energy confound      & CNN & Grad-CAM & (validity) \\
Variant robustness   & CNN & Grad-CAM vs.\ ++ & (validity) \\
Frequency profile    & SVM, RF & SHAP & frequency \\
Cross-model agreement& SVM, RF & SHAP & frequency \\
Cardiac-cycle enrichment & CNN & Grad-CAM & time \\
Outcome stratification   & CNN & Grad-CAM & time \\
\bottomrule
\end{tabular}
\end{table}

The specific estimators: for SHAP, the exact tree estimator on the random
forests and the model-agnostic kernel estimator on the
SVMs~\cite{lundberg2017unified,lundberg2020local}; for the CNN, Grad-CAM at the
last convolutional block~\cite{selvaraju2017grad}. Table~\ref{tab:xai-map}
summarises which model and method each experiment below uses and which axis it
is allowed to speak to.

\subsection{Sanity checks: is the CNN's Grad-CAM meaningful at all?}

\emph{Model: CNN. Method: Grad-CAM. Tested: whether the attribution reflects
learned weights rather than architecture or loudness.} These checks are reported
before any finding, because an attribution map means nothing until it is shown
to depend on the trained weights~\cite{adebayo2018sanity}.

\textbf{Weight randomisation.} Recomputing Grad-CAM from a randomly
re-initialised copy of the same architecture gives maps that correlate with the
trained model's at only \num{0.051} over 100 windows. The maps therefore reflect
what the network learned, not its wiring. \emph{Pass.}

\textbf{Energy confound.} Correlating the CNN's temporal attribution profile
with the signal envelope gives \num{0.068} over 499 windows, with no window
above \num{0.5}. Attribution is not merely tracking loudness --- which matters
because S1 and S2 are the loud events, so a loudness-driven map would confound
the timing claim of Section~\ref{sec:xai-time}. \emph{Pass.}

\textbf{Variant robustness.} Grad-CAM and Grad-CAM++~\cite{chattopadhay2018grad}
maps correlate at only \num{0.392}. The temporal conclusions below are therefore
specific to the Grad-CAM configuration and should not be assumed to transfer to
other CAM variants. \emph{Reported as a limitation.}

\subsection{Where in frequency? (SHAP, classical models)}
\label{sec:xai-freq}

\emph{Models: SVM and RF, on MFCC and PWP. Method: SHAP. Tested: which
frequencies drive the classical models, and whether independent explainers
agree on the answer.} This is the only place the frequency axis is claimed, and
Grad-CAM contributes nothing to it.

\begin{figure}[t]
\centering
\includegraphics[width=\columnwidth]{fig_frequency_attribution}
\caption{SHAP attribution mapped to frequency for the four classical
feature/model pairings. PWP features map onto bands exactly; the MFCC profile is
obtained by redistributing attribution through the magnitude of the DCT basis
and is approximate.}
\label{fig:freqattr}
\end{figure}

The frequency mapping is exact for PWP and approximate for MFCC, and the
distinction governs which number to trust. A PWP feature belongs to exactly one
frequency band by construction, so its SHAP value maps onto frequency directly.
An MFCC coefficient is a weighted sum over all mel bands, so its attribution must
be redistributed through the magnitude of the DCT basis; that projection is
sign-agnostic and therefore approximate. Where the two disagree, the PWP result
is the reliable one.

\textbf{Cross-model agreement is weak.} All four classical pairings peak at the
low band edge, \SI{25}{\hertz} (Fig.~\ref{fig:freqattr}), but the mean pairwise
correlation between the four frequency profiles is only \num{0.514}, and that
average hides a wide range. The two MFCC models agree almost perfectly with each
other (\num{0.961}), whereas RF--PWP and SVM--MFCC agree at \num{0.280}. The high
figure is a statement about a shared \emph{representation}, not a shared
physiological finding: models built on the same features attribute similarly,
models built on different features do not. Agreement across independent
explainers would have been the strongest single piece of interpretability
evidence available --- it is what a genuine cross-method check looks like on the
frequency axis --- and it did not materialise. We report frequency agreement as
weak.

\textbf{The individual PWP features are plausible but not decisive.} The top two
by mean absolute SHAP value are the Shannon entropy of band~0
(25--\SI{51}{\hertz}, \SI{9.5}{\percent} of total attribution) and of band~9
(291--\SI{326}{\hertz}, \SI{7.9}{\percent}), followed by relative energy in
77--\SI{105}{\hertz}. Spectral entropy in a low band and a high band is a
sensible murmur descriptor, since murmurs are noise-like where S1 and S2 are
tonal. But their sign consistency is only \num{0.57} and \num{0.55}: they push
toward ``abnormal'' barely more often than toward ``normal''. We therefore
describe them as the features the model relies on most, not as evidence of a
murmur detector. For the MFCC models the leading features are $\mathrm{MFCC}_3$
mean and $\mathrm{MFCC}_0$ statistics; the prominence of the zeroth coefficient,
which encodes overall energy, is not encouraging given the site dependence of
Section~\ref{sec:persite}.

\subsection{Where in the cardiac cycle? (Grad-CAM, CNN)}
\label{sec:xai-time}

\emph{Model: CNN. Method: Grad-CAM. Tested: whether the CNN's attention
concentrates in systole, where clinically important adult murmurs occur.} This
is the only place the time axis is claimed, and SHAP contributes nothing to it.

\begin{table}[t]
\centering
\caption{Cardiac-cycle alignment of model attributions. $E_s$ is the enrichment (attribution mass divided by time fraction); $E_s = 1$ corresponds to uniform temporal attention.}
\label{tab:alignment}
\begin{tabular}{lrrrrrrrr}
\toprule
model & mass S1 & E S1 & mass systole & E systole & mass S2 & E S2 & mass diastole & E diastole \\
\midrule
cnn\_logmel\_gradcam & 0.186 & 0.953 & 0.221 & 0.936 & 0.158 & 0.766 & 0.435 & 1.092 \\
\bottomrule
\end{tabular}
\begin{flushleft}\footnotesize Computed on the 66\textbackslash\{\}\% of test windows the segmenter labelled with sufficient confidence.\end{flushleft}
\end{table}


\begin{figure}[t]
\centering
\includegraphics[width=\columnwidth]{fig_alignment_cnn}
\caption{Distribution of per-window systolic enrichment $E_\text{systole}$ for
the CNN. The dashed line at $E=1$ is what uniform temporal attention produces;
the distribution sits slightly to its right.}
\label{fig:alignment}
\end{figure}

An envelope-based segmenter, operating independently of the CNN, labels each
window's timeline into S1, systole, S2 and diastole. Given a unit-mass Grad-CAM
map, the enrichment for cardiac state $s$ is
\begin{equation}
E_s = \frac{\text{fraction of attribution mass inside } s}
           {\text{fraction of time occupied by } s},
\end{equation}
so $E_s = 1$ is exactly what uniform temporal attention produces, $E_s > 1$ is a
genuine preference, and $E_s < 1$ is avoidance. The normalisation is the point:
``\SI{25}{\percent} of attention fell in systole'' is uninterpretable until one
knows systole occupies \SI{23}{\percent} of the window. Because most clinically
important adult murmurs are systolic, a CNN that has learned pathology should
show $E_\text{systole} > 1$ --- a prediction fixed before the measurement.

\textbf{Confirmed in sign, negligible in magnitude.} Table~\ref{tab:alignment}
gives $E_\text{systole} = \num{1.056}$, above the uniform null by about
\SI{6}{\percent}. Against a permutation null --- the same attribution profile
shuffled in time, so its sparsity and dynamic range survive but any relation to
the cardiac cycle is destroyed --- only \SI{33}{\percent} of the 200 tested
windows reach $\alpha = 0.05$, the mean $p$-value is \num{0.43}, and the effect
size is Cohen's $d = \num{0.131}$ (below the conventional threshold for a small
effect). Fig.~\ref{fig:alignment} shows a distribution shifted right of $E = 1$
but overlapping it heavily. This is a tendency, not a property of the model.

Three conditions attach to these numbers. First, the segmenter is a simplified
envelope-based labeller, not the reference hidden semi-Markov
model~\cite{springer2016logistic}; it is used only at evaluation and never in
training, so a noisier segmenter adds variance and makes the test \emph{harder}
to pass rather than easier, and cannot manufacture a positive result. Second,
low-confidence windows are excluded, so the honest denominator for every figure
here is the \SI{66.5}{\percent} inclusion rate (2235 of 3362 windows, mean
segmenter confidence \num{0.559}, mean estimated heart rate \SI{82}{bpm}).
Third, as stated at the outset, Grad-CAM's coarse frequency axis means no
spectral claim is drawn from it; the frequency evidence is entirely SHAP's.

\subsection{The check that decides it: right for the right reason? (Grad-CAM, CNN)}

\emph{Model: CNN. Method: Grad-CAM. Tested: whether the systolic tendency
distinguishes correct from incorrect decisions --- the difference between the
CNN \emph{attending} to systole and the CNN \emph{using} systolic evidence.}

\begin{figure}[t]
\centering
\includegraphics[width=\columnwidth]{fig_gradcam_example}
\caption{Grad-CAM on a correctly classified abnormal window, with the segmenter
state track below. Maps such as this motivate the enrichment test but cannot,
on their own, support it.}
\label{fig:gradcam}
\end{figure}

If the systolic tendency reflected murmur detection, it should be stronger when
the CNN is right than when it is wrong. Stratifying the same Grad-CAM enrichment
by outcome, it is not: mean $E_\text{systole}$ is \num{1.085} on true positives
but \num{1.115} on false positives and \num{1.157} on false negatives
($n = 26$, so the last is a tendency at best). The true-minus-false-positive
difference is $\num{-0.030}$ --- the systolic focus is marginally
\emph{stronger} when the CNN raises a false alarm.

This is the clearest negative result in the project, and it bounds what the rest
of the section may claim. The maps are not artifacts --- the sanity checks
establish that --- and they do lean toward systole. But the lean does not
separate correct from incorrect decisions, so it cannot be evidence that the CNN
uses systolic information to decide. It behaves more like a fixed positional
habit of the architecture than like a learned mechanism. Fig.~\ref{fig:gradcam}
shows a representative correctly classified abnormal window: visual inspection of
such a map is what suggested the hypothesis, and the quantitative test is what
stops us concluding it.
